\documentclass[a4paper,10pt]{article}
\usepackage{multicol}
\usepackage{amsmath}
\usepackage{amssymb}
\usepackage{bm}
\usepackage{soul}
\usepackage{xcolor}
\sethlcolor{cyan}
\usepackage[margin=0.5in]{geometry}
\usepackage{graphicx}
\usepackage{times}
\renewcommand{\baselinestretch}{0.9} 

\title{Other Notes}
\author{}
\date{}

\begin{document}

\maketitle


\section{Logistic Regression}

Logistic regression is a generalized linear model used for binary outcomes.

\subsection{Model}

Let
\[
Y \in \{0,1\}, \quad X \in \mathbb{R}^p
\]

We model

\[
P(Y=1|X) = p(X)
\]

Instead of modeling the probability directly, logistic regression models the log-odds:

\[
\log \frac{P(Y=1|X)}{1-P(Y=1|X)}
=
\beta_0 + X^T\beta
\]

This is called the \textbf{logit link function}.

\subsection{Logistic Function}

Solving for the probability gives

\[
P(Y=1|X)
=
\frac{1}{1+\exp(-(\beta_0 + X^T\beta))}
\]

Let

\[
\eta = \beta_0 + X^T\beta
\]

then

\[
p(X) = \frac{e^\eta}{1+e^\eta}.
\]

\subsection{Likelihood}

Assume independent observations

\[
(Y_i, X_i), \quad i=1,\dots,n
\]

with

\[
Y_i \sim \text{Bernoulli}(p_i)
\]

where

\[
p_i = P(Y_i=1|X_i).
\]

The likelihood is

\[
L(\beta)
=
\prod_{i=1}^n
p_i^{Y_i}(1-p_i)^{1-Y_i}.
\]

\subsection{Log-Likelihood}

\[
\ell(\beta)
=
\sum_{i=1}^n
\left[
Y_i \log p_i
+
(1-Y_i)\log(1-p_i)
\right].
\]

Substituting the logistic form gives

\[
\ell(\beta)
=
\sum_{i=1}^n
\left[
Y_i X_i^T\beta
-
\log(1+\exp(X_i^T\beta))
\right].
\]

\subsection{Estimation}

Parameters are estimated using maximum likelihood estimation (MLE).  
There is no closed-form solution, so numerical methods such as Newton–Raphson or IRLS are used.

The score function is

\[
U(\beta)
=
\sum_{i=1}^n
X_i(Y_i - p_i).
\]

The Fisher information matrix is

\[
I(\beta)
=
\sum_{i=1}^n
p_i(1-p_i) X_i X_i^T.
\]

Asymptotically,

\[
\hat{\beta}
\sim
N(\beta, I(\beta)^{-1}).
\]

\subsection{Interpretation}

Each coefficient represents a log-odds change:

\[
\beta_j
=
\log(\text{odds ratio})
\]

and

\[
\exp(\beta_j)
\]

is the multiplicative change in odds when $x_j$ increases by one unit.

\subsection{Classification}

Predicted probability:

\[
\hat p(X)
\]

Classification rule:

\[
\hat Y =
\begin{cases}
1 & \hat p(X) > c \\
0 & \text{otherwise}
\end{cases}
\]

where $c$ is a threshold (often $0.5$).

\subsection{Regularization}

To avoid overfitting:

Ridge logistic regression:

\[
\ell(\beta) - \lambda \|\beta\|_2^2
\]

Lasso logistic regression:

\[
\ell(\beta) - \lambda \|\beta\|_1.
\]

\section{EM Algorithm for MLE with Missing Data}

Suppose the complete data are
\[
X = (Y,Z)
\]

where

\begin{itemize}
\item $Y$ : observed data
\item $Z$ : missing (latent) data
\end{itemize}

Assume the model

\[
p(Y,Z|\theta).
\]

\subsection{Observed Likelihood}

The observed likelihood is

\[
L(\theta)
=
p(Y|\theta)
=
\sum_Z p(Y,Z|\theta)
\]

(or an integral if $Z$ is continuous).

The log-likelihood is

\[
\ell(\theta)
=
\log p(Y|\theta)
=
\log \sum_Z p(Y,Z|\theta).
\]

This is often difficult to maximize directly.

\subsection{EM Algorithm}

The EM algorithm maximizes the likelihood iteratively.

Define

\[
Q(\theta|\theta^{(t)})
=
E_{Z|Y,\theta^{(t)}}
[
\log p(Y,Z|\theta)
].
\]

\subsection{Algorithm Steps}

\textbf{E-step}

\[
Q(\theta|\theta^{(t)})
=
E_{Z|Y,\theta^{(t)}}[
\log p(Y,Z|\theta)
].
\]

Compute the expected complete-data log-likelihood.

\textbf{M-step}

\[
\theta^{(t+1)}
=
\arg\max_\theta Q(\theta|\theta^{(t)}).
\]

Repeat until convergence.

The EM algorithm guarantees

\[
\ell(\theta^{(t+1)}) \ge \ell(\theta^{(t)}).
\]

\subsection{Example: Gaussian Mixture Model}

Assume

\[
Y_i \sim \pi N(\mu_1,\sigma^2) + (1-\pi)N(\mu_2,\sigma^2).
\]

Introduce latent variable

\[
Z_i =
\begin{cases}
1 & \text{component 1} \\
0 & \text{component 2}.
\end{cases}
\]

\subsection{E-step}

Compute posterior probability

\[
\gamma_i
=
P(Z_i=1|Y_i,\theta^{(t)})
=
\frac{
\pi^{(t)} f_1(Y_i)
}{
\pi^{(t)} f_1(Y_i) + (1-\pi^{(t)}) f_2(Y_i)
}.
\]

\subsection{M-step}

Update parameters

\[
\pi^{(t+1)}
=
\frac{1}{n}\sum_i \gamma_i.
\]

\[
\mu_1^{(t+1)}
=
\frac{\sum_i \gamma_i Y_i}{\sum_i \gamma_i}.
\]

\[
\mu_2^{(t+1)}
=
\frac{\sum_i (1-\gamma_i)Y_i}{\sum_i (1-\gamma_i)}.
\]

\[
\sigma^2
=
\frac{
\sum_i \gamma_i (Y_i-\mu_1)^2
+
\sum_i (1-\gamma_i)(Y_i-\mu_2)^2
}{n}.
\]

\section{Augmented Inverse Probability Weighting (AIPW) / Doubly Robust Estimation}

\subsection{Setup}

Let

\begin{itemize}
\item $A \in \{0,1\}$ : treatment
\item $X$ : covariates
\item $Y$ : observed outcome
\end{itemize}

Potential outcomes:

\[
Y(1), Y(0)
\]

Observed outcome:

\[
Y = AY(1) + (1-A)Y(0)
\]

The target estimand is the Average Treatment Effect (ATE)

\[
\tau = E[Y(1) - Y(0)].
\]

\subsection{Identification Assumptions}

\textbf{Consistency}

\[
Y = Y(A)
\]

\textbf{Unconfoundedness}

\[
(Y(1),Y(0)) \perp A \mid X
\]

\textbf{Positivity}

\[
0 < P(A=1|X) < 1
\]

Under these assumptions

\[
\tau =
E_X
\left[
E[Y|A=1,X] - E[Y|A=0,X]
\right].
\]

\subsection{Outcome Regression}

Define

\[
\mu_a(X) = E[Y|A=a,X].
\]

Estimator

\[
\hat\tau_{OR}
=
\frac{1}{n}
\sum_i
\left(
\hat\mu_1(X_i) - \hat\mu_0(X_i)
\right).
\]

\subsection{Inverse Probability Weighting}

Define the propensity score

\[
e(X) = P(A=1|X).
\]

Estimator

\[
\hat\tau_{IPW}
=
\frac{1}{n}
\sum_i
\left(
\frac{A_i Y_i}{\hat e(X_i)}
-
\frac{(1-A_i)Y_i}{1-\hat e(X_i)}
\right).
\]

\subsection{Augmented IPW (Doubly Robust Estimator)}

The AIPW estimator is

\[
\hat\tau_{DR}
=
\frac{1}{n}
\sum_i
\Bigg[
\hat\mu_1(X_i) - \hat\mu_0(X_i)
+
\frac{A_i}{\hat e(X_i)}(Y_i-\hat\mu_1(X_i))
-
\frac{1-A_i}{1-\hat e(X_i)}(Y_i-\hat\mu_0(X_i))
\Bigg].
\]

\subsection{Double Robustness}

The estimator is consistent if either

\begin{itemize}
\item the outcome model $\mu_a(X)$ is correctly specified, or
\item the propensity score model $e(X)$ is correctly specified.
\end{itemize}

\subsection{Efficient Influence Function}

The efficient influence function for the ATE is

\[
\phi(O)
=
\mu_1(X) - \mu_0(X)
+
\frac{A}{e(X)}(Y-\mu_1(X))
-
\frac{1-A}{1-e(X)}(Y-\mu_0(X))
-
\tau.
\]

The AIPW estimator is obtained by solving

\[
\frac{1}{n}\sum_i \phi(O_i) = 0.
\]

\section{Kernel Regression}

Kernel regression is a nonparametric method for estimating the regression function

\[
m(x)=E[Y\mid X=x].
\]

Suppose we observe

\[
(X_i,Y_i), \quad i=1,\dots,n.
\]

The main idea is to estimate $m(x)$ by taking a locally weighted average of the responses $Y_i$, where observations with $X_i$ close to $x$ receive larger weights.

\subsection{Kernel Function}

A kernel function $K(\cdot)$ is a nonnegative weighting function, usually symmetric and integrating to 1.

Examples include:

\paragraph{Gaussian kernel}
\[
K(u)=\frac{1}{\sqrt{2\pi}}e^{-u^2/2}
\]

\paragraph{Uniform kernel}
\[
K(u)=\frac{1}{2}\mathbf{1}(|u|\le 1)
\]

\paragraph{Epanechnikov kernel}
\[
K(u)=\frac{3}{4}(1-u^2)\mathbf{1}(|u|\le 1)
\]

\subsection{Bandwidth}

The bandwidth $h>0$ controls the amount of smoothing.

\begin{itemize}
\item Small $h$: less smoothing, lower bias, higher variance
\item Large $h$: more smoothing, higher bias, lower variance
\end{itemize}

\subsection{Nadaraya--Watson Estimator}

The kernel regression estimator is

\[
\hat m(x)
=
\frac{\sum_{i=1}^n K\left(\frac{x-X_i}{h}\right)Y_i}
{\sum_{i=1}^n K\left(\frac{x-X_i}{h}\right)}.
\]

Equivalently, define the weights

\[
w_i(x)=
\frac{K\left(\frac{x-X_i}{h}\right)}
{\sum_{j=1}^n K\left(\frac{x-X_j}{h}\right)},
\]

then

\[
\hat m(x)=\sum_{i=1}^n w_i(x)Y_i.
\]

Thus $\hat m(x)$ is a weighted average of the observed responses near $x$.

\subsection{Local Constant Interpretation}

Kernel regression can be viewed as a local constant fit. For a fixed target point $x$, consider

\[
\sum_{i=1}^n K\left(\frac{x-X_i}{h}\right)(Y_i-c)^2.
\]

The minimizer over $c$ is exactly the Nadaraya--Watson estimator.

\subsection{Bias-Variance Tradeoff}

Kernel regression has the usual bias-variance tradeoff:

\begin{itemize}
\item smaller bandwidth gives lower bias but higher variance
\item larger bandwidth gives higher bias but lower variance
\end{itemize}

Under smoothness assumptions, in one dimension the bias is typically of order $h^2$ and the variance is of order $1/(nh)$.

\subsection{Boundary Effect}

Kernel regression often has larger bias near the boundary of the covariate range, because fewer observations are available on one side of the target point.

\subsection{Extensions}

For multivariate covariates $X\in\mathbb{R}^p$, kernel regression extends naturally, but suffers from the curse of dimensionality.

A common improvement is local linear regression, which reduces boundary bias.

\subsection{Summary}

Kernel regression is a simple and important nonparametric regression method that estimates the conditional mean function by locally weighted averaging. Its performance depends strongly on the choice of bandwidth.

\section{Linear Algebra}
\subsection{Derivation}
\[
\begin{aligned}
\frac{\partial}{\partial x}(a^T x) &= a \\
\frac{\partial}{\partial x}(x^T A x) &= (A+A^T)x \\
\frac{\partial}{\partial x}(x^T A x) &= 2Ax,\quad A=A^T \\
\frac{\partial}{\partial x}\|Ax-b\|^2 &= 2A^T(Ax-b) \\
\frac{\partial}{\partial x}(A+Bx)^T(A+Bx) &= 2B^T(A+Bx)
\end{aligned}
\] 
\subsection{Covariance matrix}
\[
\mathrm{Cov}(aX)=a^2\Sigma
\]

\[
\mathrm{Cov}(AX)=A\Sigma A^T
\]

\[
\mathrm{Cov}(A+BX)=B\Sigma B^T
\]

\subsection{Trace}
\[
\operatorname{tr}(A)=\sum_i a_{ii}
\]

\[
\operatorname{tr}(A+B)
=
\operatorname{tr}(A)+\operatorname{tr}(B)
\]

\[
\operatorname{tr}(AB)=\operatorname{tr}(BA)
\]

\[
\operatorname{tr}(ABC)
=
\operatorname{tr}(BCA)
=
\operatorname{tr}(CAB)
\]

\[
\operatorname{tr}(P^{-1}AP)=\operatorname{tr}(A)
\]

\[
\operatorname{tr}(A)=\sum_i \lambda_i
\]

\[
\operatorname{tr}(A^k)=\sum_i \lambda_i^k
\]

\[
\operatorname{tr}(A^T A)=\lVert A\rVert_F^2
\]

\[
x^T A x=\operatorname{tr}(Axx^T)
\]

\[
P^2=P
\quad\Longrightarrow\quad
\operatorname{tr}(P)=\operatorname{rank}(P)
\]

\[
\operatorname{tr}(AB-BA)=0
\]

\[
\det(e^A)=e^{\operatorname{tr}(A)}
\]

\section{Linear Regression: Important Concepts and Conclusions}

\subsection{Model}

For observations $i=1,\dots,n$, the linear regression model is

\[
Y_i = \beta_0 + \beta_1 X_{i1} + \cdots + \beta_p X_{ip} + \varepsilon_i.
\]

In matrix form,

\[
Y = X\beta + \varepsilon.
\]

Here $Y$ is the response vector, $X$ is the design matrix, $\beta$ is the coefficient vector, and $\varepsilon$ is the error vector.

The model is called \emph{linear} because it is linear in the coefficients $\beta$.

\subsection{Conditional Mean Interpretation}

The regression model describes the conditional mean

\[
E[Y\mid X] = X\beta.
\]

In multiple regression, $\beta_j$ represents the change in the conditional mean of $Y$ associated with a one-unit increase in $X_j$, holding the other predictors fixed.

\subsection{Ordinary Least Squares (OLS)}

The ordinary least squares estimator minimizes the residual sum of squares

\[
RSS(\beta) = (Y-X\beta)^T(Y-X\beta).
\]

The OLS estimator is

\[
\hat\beta = (X^T X)^{-1}X^T Y,
\]

provided $X^TX$ is invertible.

The normal equations are

\[
X^T(Y-X\hat\beta)=0.
\]

\subsection{Fitted Values and Residuals}

The fitted values are

\[
\hat Y = X\hat\beta,
\]

and the residuals are

\[
e = Y-\hat Y.
\]

Define the hat matrix

\[
H = X(X^TX)^{-1}X^T.
\]

Then

\[
\hat Y = HY, \qquad e = (I-H)Y.
\]

Important facts:

\begin{itemize}
\item $H$ is symmetric and idempotent
\item $\hat Y$ is the projection of $Y$ onto the column space of $X$
\item residuals are orthogonal to the column space of $X$
\end{itemize}

\subsection{Assumptions}

A key assumption is

\[
E[\varepsilon \mid X] = 0,
\]

which implies

\[
E[Y\mid X] = X\beta.
\]

A common variance assumption is

\[
\mathrm{Var}(\varepsilon \mid X)=\sigma^2 I.
\]

Sometimes we also assume normality:

\[
\varepsilon \mid X \sim N(0,\sigma^2 I).
\]

\subsection{Main Properties of OLS}

If $E[\varepsilon\mid X]=0$, then OLS is unbiased:

\[
E[\hat\beta \mid X]=\beta.
\]

If in addition $\mathrm{Var}(\varepsilon\mid X)=\sigma^2 I$, then

\[
\mathrm{Var}(\hat\beta \mid X)=\sigma^2 (X^TX)^{-1}.
\]

Under these assumptions, OLS is the BLUE (Best Linear Unbiased Estimator), by the Gauss--Markov theorem.

If we also assume normality, then

\[
\hat\beta \mid X \sim N\bigl(\beta,\sigma^2(X^TX)^{-1}\bigr).
\]

\subsection{Estimating the Error Variance}

The residual sum of squares is

\[
RSS = \sum_{i=1}^n e_i^2.
\]

An unbiased estimator of $\sigma^2$ is

\[
\hat\sigma^2 = \frac{RSS}{n-p-1}.
\]

\subsection{Inference for Coefficients}

The standard error of $\hat\beta_j$ is

\[
SE(\hat\beta_j)
=
\sqrt{\hat\sigma^2\bigl[(X^TX)^{-1}\bigr]_{jj}}.
\]

To test

\[
H_0:\beta_j=0,
\]

use the statistic

\[
t=\frac{\hat\beta_j}{SE(\hat\beta_j)},
\]

which follows a $t$ distribution with $n-p-1$ degrees of freedom under the normal model.

A confidence interval is

\[
\hat\beta_j \pm t_{n-p-1,1-\alpha/2}SE(\hat\beta_j).
\]

\subsection{Overall Model Test}

To test

\[
H_0:\beta_1=\cdots=\beta_p=0,
\]

use the $F$-statistic

\[
F=\frac{(SSR/p)}{(RSS/(n-p-1))}.
\]

\subsection{ANOVA Decomposition and $R^2$}

The total variation decomposes as

\[
TSS = SSR + RSS.
\]

The coefficient of determination is

\[
R^2 = \frac{SSR}{TSS} = 1-\frac{RSS}{TSS}.
\]

Adjusted $R^2$ is

\[
R^2_{adj}
=
1-\frac{RSS/(n-p-1)}{TSS/(n-1)}.
\]

\subsection{Prediction}

For a new covariate vector $x_0$, the predicted mean response is

\[
\hat Y_0 = x_0^T\hat\beta.
\]

A confidence interval for the mean response is narrower than a prediction interval for a new observation, because prediction intervals must also include the new noise term.

\subsection{Important Practical Issues}

\paragraph{Multicollinearity}
Highly correlated predictors can make $(X^TX)^{-1}$ unstable, leading to large standard errors.

\paragraph{Omitted Variable Bias}
If an omitted variable is correlated with included predictors, then $E[\varepsilon\mid X]\neq 0$, and OLS becomes biased.

\paragraph{Heteroscedasticity}
If $\mathrm{Var}(\varepsilon\mid X)$ is not constant, OLS can still be unbiased, but the usual variance formula is incorrect.

\paragraph{Influential Observations}
Some observations may have unusually large influence on the fitted model, so diagnostic checks are important.

\subsection{Summary}

Linear regression is built around three core ideas:

\begin{itemize}
\item modeling the conditional mean: $E[Y\mid X]=X\beta$
\item estimating coefficients by least squares: $\hat\beta=(X^TX)^{-1}X^TY$
\item doing inference using assumptions on the error distribution
\end{itemize}

\section{Distribution}
\subsection{Multinomial Distribution}
$$
(X_1,\dots,X_k)
\sim
\operatorname{Multinomial}(n;p_1,\dots,p_k)
$$

$$
P(X_1=x_1,\dots,X_k=x_k)
=
\frac{n!}{\prod_i x_i!}
\prod_i p_i^{x_i}
$$

$$
E[X_i]=np_i
$$

$$
\operatorname{Var}(X_i)=np_i(1-p_i),
\qquad
\operatorname{Cov}(X_i,X_j)=-np_ip_j
$$

$$
\sum_{i=1}^k X_i=n
$$
\end{document}
